\section{Teoría de robótica móvil: ruedas, ICR y maniobrabilidad}
\label{sec:robotica}

El modelo de \cref{sec:modelo-fisico} parte de la restricción de Pfaff de la rueda diferencial. Esta sección da el marco general del que esa restricción procede \cite{siegwart2011introduction} y lo que ese marco impone a
una coalición, que resulta ser un teorema geométrico.

\subsection{Restricciones de rueda}

Sea $\xi=(x,y,\theta)$ la pose del chasis y $R(\theta)$ la rotación al marco del
robot. Una rueda estándar situada a distancia $\ell$ del centro del chasis, con
ángulo polar $\alpha$, orientación $\beta$ respecto del radio, radio $r$ y
velocidad angular $\dot\varphi$, impone dos restricciones:
\begin{align}
  \text{rodadura:}&\quad
  \bigl(\sin(\alpha+\beta),\ -\cos(\alpha+\beta),\ -\ell\cos\beta\bigr)\,R(\theta)\dot\xi
  \;=\;r\dot\varphi, \label{eq:rodadura}\\
  \text{no deslizamiento:}&\quad
  \bigl(\cos(\alpha+\beta),\ \sin(\alpha+\beta),\ \ell\sin\beta\bigr)\,R(\theta)\dot\xi
  \;=\;0. \label{eq:deslizamiento}
\end{align}
La rueda \emph{fija} tiene $\beta$ constante; la \emph{orientable} tiene
$\beta(t)$ actuada; la \emph{castor} añade un descentramiento que hace su
restricción de deslizamiento no acotante; las ruedas \emph{suecas} y
\emph{esféricas} no imponen \eqref{eq:deslizamiento} en absoluto. Apilando las
restricciones de deslizamiento de las $N_f$ ruedas fijas y $N_s$ orientables se
obtiene $C_1(\beta_s)R(\theta)\dot\xi=0$, y el conjunto de velocidades
admisibles del chasis es $\ker C_1(\beta_s)$.

\begin{definicion}[Movilidad, direccionabilidad y maniobrabilidad]
\label{def:maniobrabilidad}
$\delta_m=3-\operatorname{rank}C_1(\beta_s)$ es el grado de \emph{movilidad}:
las velocidades que el chasis puede imprimir instantáneamente sin cambiar la
orientación de sus ruedas. $\delta_s=\operatorname{rank}C_{1s}(\beta_s)$ es el
grado de \emph{direccionabilidad}: las orientaciones de rueda que puede cambiar.
La \emph{maniobrabilidad} es $\delta_M=\delta_m+\delta_s$.
\end{definicion}

\begin{center}\small
\begin{tabular}{@{}lccc l@{}}
\toprule
Chasis & $\delta_m$ & $\delta_s$ & $\delta_M$ & holonómico \\
\midrule
omnidireccional (tres ruedas suecas) & 3 & 0 & 3 & sí \\
diferencial (dos fijas + castor) & 2 & 0 & 2 & no \\
omni-direccionable (tres orientables) & 1 & 2 & 3 & no \\
triciclo / Ackermann & 1 & 1 & 2 & no \\
dos direccionables (\emph{two-steer}) & 1 & 2 & 3 & no \\
\bottomrule
\end{tabular}
\end{center}

Un robot es holonómico si sus grados de libertad diferenciables igualan a los del
espacio de trabajo, $\delta_m=3$; solo la primera fila lo es. El AMR de este
trabajo es la segunda: $\delta_m=2$, $\delta_s=0$, y ese es el origen de
\eqref{eq:pfaff}. La maniobrabilidad $\delta_M$ dice qué poses puede alcanzar;
la movilidad $\delta_m$ dice cómo puede moverse ahora. Que $\delta_M=\delta_m$
en el diferencial significa que no hay nada que reorientar: toda maniobra es un
cambio de velocidades de rueda.

\subsection{El centro instantáneo de rotación}

Cada rueda estándar define una \emph{línea de movimiento nulo}: la perpendicular
a su plano por su centro. En todo instante el chasis rota alrededor de un punto,
el \emph{centro instantáneo de rotación} (ICR), que debe estar sobre todas esas
líneas a la vez; si están en posición general y son más de una, el ICR está
determinado y $\delta_m$ cae. Es una construcción sobre el número de
restricciones, no sobre el número de ruedas: una bicicleta y un Ackermann tienen
$\delta_m=1$ con dos y cuatro ruedas. El diferencial tiene sus dos ruedas sobre
el mismo eje, luego una sola línea, y el ICR es libre sobre ella: eso es
$\delta_m=2$. Con $\zeta=(V,\Omega)$ el \emph{twist} del chasis en su propio
marco, el ICR está en $r_{\mathrm{ICR}}=SV/\Omega$, a distancia
$\lVert V\rVert/|\Omega|$ sobre la perpendicular a $V$.

\subsection{Lo que el ICR impone a una coalición}

Una coalición acoplada rígidamente a la carga es un solo cuerpo, y un cuerpo
tiene un solo ICR. Cada AMR miembro aporta su línea de movimiento nulo. Esto da
la lectura geométrica de \cref{pr:rigidez}.

\begin{teorema}[ICR común]
\label{th:icr-comun}
Sea $\zeta_\ell=(V_\ell,\Omega_\ell)$ el \emph{twist} de la carga con
$\Omega_\ell\neq0$, y $r_{\mathrm{ICR}}=SV_\ell/\Omega_\ell$. Para el AMR $i$ con
rumbo $e_i$ y centro de eje en $r_i$ (marco de la carga),
\[
  [A_{\mathrm{kin}}]_i\,\zeta_\ell\;=\;\Omega_\ell\,e_i^{\!\top}\bigl(r_i-r_{\mathrm{ICR}}\bigr).
\]
En consecuencia, la fila $i$ de la restricción de rigidez se anula si y solo si
el ICR de la carga está sobre la línea de movimiento nulo de la rueda $i$, y la
coalición puede ejecutar $\zeta_\ell$ sin deslizamiento lateral si y solo si
\emph{todas} sus líneas de movimiento nulo concurren en $r_{\mathrm{ICR}}$. Con
$\Omega_\ell=0$ la condición es que todas sean paralelas a $V_\ell^{\perp}$.
\end{teorema}

\begin{proof}
La velocidad del punto $r_i$ es $V_\ell+\Omega_\ell Sr_i=\Omega_\ell S(r_i-r_{\mathrm{ICR}})$
usando $V_\ell=-\Omega_\ell S\,r_{\mathrm{ICR}}$ y $S^{2}=-I$. Su componente
lateral es $e_i^{\perp\top}\Omega_\ell S(r_i-r_{\mathrm{ICR}})$, y como
$e_i^{\perp}=Se_i$ y $S^{\!\top}S=I$, vale $\Omega_\ell\,e_i^{\!\top}(r_i-r_{\mathrm{ICR}})$,
que es la fila de \cref{pr:rigidez} evaluada en $\zeta_\ell$. Se anula si y solo
si $r_i-r_{\mathrm{ICR}}\perp e_i$, es decir, si el ICR está sobre la
perpendicular al rumbo por $r_i$.
\end{proof}

\begin{observacion}[Verificación y consecuencias]
\label{obs:icr-comun}
La identidad se comprueba sobre dos mil configuraciones aleatorias sin una sola
discrepancia por encima de $10^{-9}$. Tres consecuencias siguen.

Primera, la dimensión $3-\operatorname{rank}A_{\mathrm{kin}}$ de
\cref{pr:ruta-ruedas} es el número de parámetros libres del ICR compatible con
todas las líneas: $2$ si todas coinciden (un solo eje común), $1$ si son
paralelas distintas (ICR en el infinito, solo traslación), $0$ si concurren en
un punto (ese ICR y solo ese), y la coalición está bloqueada si no concurren.

Segunda, con contactos fuera del eje (\cref{obs:rigidez-alcance}) la línea de
movimiento nulo sigue existiendo —es propiedad de la rueda, no del contacto—,
pero el punto de contacto adelantado puede tener velocidad lateral, de modo que
la concurrencia deja de ser obligatoria y pasa a ser económica: cuanto más lejos
esté el ICR de la línea, mayor es la $\omega_i$ que \eqref{eq:N-eje} exige.

Tercera, y es una regla de formación: para que una coalición diferencial gire
sobre sí misma ($V_\ell=0$, ICR en el centro de la carga) todos los ejes deben
apuntar radialmente al centro; para que traslade en línea recta, todos deben ser
paralelos entre sí y perpendiculares a $V_\ell$. Ninguna formación satisface
ambas a la vez, y esa es la razón geométrica de que el reparto de puestos deba
mirar la ruta —\cref{pr:ruta-ruedas}— y no solo el \emph{wrench}.
\end{observacion}
